\section{Descriptive Analysis: The Case for a Kernel Backbone}
\label{sec:descriptive}

Before asking what exogenous information adds, we fix the instrument that will
measure it. This section argues that the target's own autoregressive structure
determines that instrument, and determines it as a \emph{truncated power-law
convolution kernel modulated by a session clock}---of which the standard HAR
regression is a coarse quadrature. The argument matters for the rest of the
paper because every attribution in Sections~\ref{sec:marginal}
and~\ref{sec:linvsnonlin} is measured as an increment over this backbone, and
an incorrectly specified backbone donates its own approximation error to
whatever is being credited.

\subsection{The Panel and the Dense-but-Weak Fact}
\label{sec:dense_weak}

The design matrix has 529 columns. Its covariance concentrates in a handful of
directions, but the forecasting signal does not live in those directions. Three
measurements establish this. They predate the frozen battery and were made on
the earlier all-buckets cache; we report them because the regime they describe
is what the rest of the paper is about, and we mark the levels that are not
panel-comparable:

\begin{itemize}
  \item \textbf{Low-rank compression destroys skill.} Principal-component
    truncation of the exogenous block retains most of the variance
    ($K=5$ recovers 72\%, $K=20$ recovers 90\%) but loses to dense shrinkage at
    every penalty setting; the best PCA-20 configuration reaches QLIKE
    $0.12890$ against $0.12530$ for the dense elastic net on the same cache.
    The signal is not in the leading eigendirections. \emph{Panel note:} these
    two levels come from the earlier all-buckets cache, not the frozen battery,
    and are comparable to each other but not to Table~\ref{tab:class}.
    \pending{PCA/PLS truncation arms on the frozen panel}
  \item \textbf{$\ell_1$ selection recovers a small unstable core plus a wide
    shallow tail.} Elastic net at its optimum retains roughly 98--120 of 529
    columns; pushing sparsity harder costs performance monotonically.
  \item \textbf{No individual column is strong.} The per-feature effects are
    uniformly small, and the strongest single engineered feature contributes
    less than the aggregate of the tail it sits in.
\end{itemize}

We use \emph{dense but weak} throughout as shorthand for this configuration:
sparse in no useful coordinate basis, not low-rank in data space, and with no
individually decisive column. Section~\ref{sec:linvsnonlin} shows what this
implies for model class; the immediate consequence is that the backbone cannot
be discovered by selection and must be specified.

\subsection{HAR Is a Quadrature Rule}
\label{sec:har_quadrature}

Write the general one-step linear forecast as a convolution of the target's own
past against a decay kernel $\varphi$:
\begin{equation}
  \hat{y}_t \;=\; \mu \;+\; \sum_{u \ge 1} \varphi(u)\, y_{t-u}.
  \label{eq:convolution}
\end{equation}
HAR does not estimate $\varphi$ freely. It regresses on trailing rolling means
(boxcars) over a ladder of rungs $r_k$:
\begin{equation}
  \hat{y}_t \;=\; \mu \;+\; \sum_{k} w_k \cdot \frac{1}{r_k}\sum_{u=1}^{r_k} y_{t-u}.
  \label{eq:har_boxcar}
\end{equation}
Collecting the total coefficient that Equation~\eqref{eq:har_boxcar} places on
an individual lag $u$ gives the \emph{effective kernel} implied by a fitted HAR:
\begin{equation}
  \boxed{\;\varphi(u) \;=\; \sum_{k \,:\, r_k \ge u} \frac{w_k}{r_k}\;}
  \label{eq:eff_kernel}
\end{equation}
which is a decreasing step function with breakpoints at the rungs. Under this
reading the $|\mathcal{L}|$ fitted HAR coefficients are not $|\mathcal{L}|$
distinct facts about volatility at $|\mathcal{L}|$ distinct horizons: they are
quadrature weights on a piecewise-constant approximation to a single underlying
curve, with the rungs as knots. Figure~\ref{fig:target_kernel} plots the fitted
boxcar coefficients and the implied $\varphi$.

\begin{figure}[H]
\centering
\treefigorpending{\textwidth}{../figs_op/fig3_target_kernel.png}
\caption{Left: self-kernel boxcar coefficients, averaged over walk-forward
refits. Right: the implied effective kernel of Equation~\eqref{eq:eff_kernel}
on log-log axes, against a fitted power law. The straight-line behaviour on
log-log axes is the object this section is about. Only the positive part of
$\varphi$ is plotted; Section~\ref{sec:kernel_caveats} discusses the negative
lobes, which are not an incidental detail.}
\label{fig:target_kernel}
\end{figure}

\subsection{The Kernel Is Power-Law and Truncated}
\label{sec:powerlaw}

The implied kernel is approximately linear on log-log axes with slope
$-\beta$, $\hat\beta \approx 1.36$ \pending{$\hat\beta$ needs a formal
estimator, a standard error, and stability across refits; it is currently a
descriptive log-log slope}. A single-parameter power law
$\varphi(u) \propto u^{-\hat\beta}$ approximately reproduces the entire fitted
ladder, which is the strongest available evidence that the ladder is
approximating one curve rather than measuring several components.

If the kernel is genuinely power-law---that is, scale-free---then the
quadrature reading of Section~\ref{sec:har_quadrature} makes three predictions
that are not made by the conventional reading of HAR as a
daily/weekly/monthly cascade. All three are testable on the shared panel, and
all three are confirmed.

\paragraph{Prediction 1: a scale-free kernel wants a scale-free grid, and
refining the grid pays.} Uniform knots in $\log$-lag are the natural quadrature
for a power law, and halving the geometric base doubles the knot density per
decade. Error should therefore fall monotonically as the base falls.
Table~\ref{tab:ladder} reports the measured ladder.

\begin{table}[H]
\centering
\caption{HAR ladder base sweep: OLS on a geometric lag ladder with base $b$,
all arms on the shared panel ($n = 218{,}934$), all differences measured
against the same incumbent. ``explicit'' is the conventional
daily/weekly/monthly HAR lag structure of \citet{Corsi2009} transposed to the
30-minute grid. Within each base we report the best cap; the cap is a
near-non-lever (within-base spread $\approx 0.0001$) while the base is not
(across-base spread $\approx 0.0014$). Negative DM $t$ favours the ladder arm.}
\label{tab:ladder}
\begin{tabular}{llrr}
\toprule
Ladder base & Arm & QLIKE & DM $t$ vs.\ incumbent \\
\midrule
$b=2$      & \texttt{b2\_c3125}       & \textbf{0.13332} & $-10.58$ \\
$b=3$      & \texttt{b3\_c3125}       & 0.13349 & $-10.19$ \\
$b=4$      & \texttt{b4\_c3125}       & 0.13361 & $-8.47$ \\
$b=5$      & \texttt{b5\_c3125} (incumbent) & 0.13415 & --- \\
$b=6$      & \texttt{b6\_c3125}       & 0.13433 & $+3.12$ \\
explicit   & \texttt{explicit\_c3125} & 0.13439 & $+3.10$ \\
$b=8$      & \texttt{b8\_c3125}       & 0.13476 & $+7.31$ \\
\bottomrule
\end{tabular}
\end{table}

The ordering is monotone in the base across all six bases, with
DM $t = -10.6$ for $b{=}2$ against the incumbent $b{=}5$. The signature is
precisely that of a quadrature error: \emph{how finely the kernel is sampled
matters, where the sum is truncated barely does}. This is difficult to explain
if the rungs were recovering genuinely distinct daily, weekly and monthly
components, and easy to explain if they are knots on one smooth curve. We note
that the conventional explicit HAR structure is not merely suboptimal but
\emph{worse than the incumbent}, at DM $t = +3.1$.

\paragraph{Prediction 2: power law, not exponential.} A single exponential
kernel fitted in place of the power law is $+0.011$ QLIKE worse---roughly an
order of magnitude larger than any effect elsewhere in this paper.

\paragraph{Prediction 3: truncation is real but shallow.} Direct power-law
convolution features degrade sharply when the convolution window is extended
too far ($U = 2048$ is catastrophic; $U = 256$ is the operating point), while
the boxcar ladder's cap is a near-non-lever. We attribute the difference to
normalisation---the direct features normalise weights to sum to one, so
extending $U$ redistributes mass away from recent lags, whereas the boxcar cap
adds a rung without disturbing the others---but this reconciliation is
currently an argument rather than a measurement.

\subsection{The Hawkes Reading}
\label{sec:hawkes}

Equation~\eqref{eq:convolution} with a power-law $\varphi$ is, in discrete
time, the conditional-intensity equation of a self-exciting point process
\citep{Hawkes1971}. This subsection makes the correspondence explicit, states
what it buys, and---at equal length---states where it does not yet hold.

\subsubsection{From point process to variance regression}

Let $N$ be a linear Hawkes process with conditional intensity
\begin{equation}
  \lambda(t) \;=\; \mu \;+\; \int_{-\infty}^{t}\varphi(t-s)\,dN(s),
  \qquad \varphi \ge 0,
  \label{eq:hawkes_intensity}
\end{equation}
and let the price be a pure-jump process driven by $N$. Realized variance over
a bar is then, up to the squared jump size, the event count over that bar, so
that $RV_t \propto N(t-\Delta, t)$ and
\begin{equation}
  \mathbb{E}\!\left[RV_t \mid \mathcal{F}_{t-1}\right]
  \;\propto\; \int_{\text{bar}} \mathbb{E}[\lambda(s)]\,ds
  \;=\; \tilde{\mu} + \sum_{u \ge 1} \bar{\varphi}(u)\, RV_{t-u},
  \label{eq:hawkes_to_rv}
\end{equation}
where $\bar\varphi$ is $\varphi$ aggregated to the bar grid. Equation
\eqref{eq:hawkes_to_rv} is Equation~\eqref{eq:convolution}. The correspondence
is therefore not an analogy: for a linear Hawkes process the conditional
intensity is exactly linear in the observed past, so a regression of counts on
past counts identifies $\varphi$ itself rather than a proxy for it. What we
estimate is a bar-aggregated Hawkes kernel, subject to the caveats of
Section~\ref{sec:kernel_caveats}.

\subsubsection{The branching ratio is already estimated}

The central quantity of Hawkes theory is the branching ratio
$n = \int_0^\infty \varphi(u)\,du$: the expected number of direct offspring per
event, equivalently the fraction of activity that is endogenously generated.
Stationarity requires $n < 1$, and $n \to 1$ is the ``critical reflexivity''
regime that has been the subject of an extended empirical debate
\citep{HardimanBercotBouchaud2013, FilimonovSornette2012}.

For the boxcar parameterisation of Equation~\eqref{eq:har_boxcar} there is an
exact identity. Substituting Equation~\eqref{eq:eff_kernel},
\begin{equation}
  \sum_{u=1}^{r_{\max}} \varphi(u)
  \;=\; \sum_{u}\sum_{k\,:\,r_k \ge u} \frac{w_k}{r_k}
  \;=\; \sum_k \frac{w_k}{r_k}\,\bigl|\{u : 1 \le u \le r_k\}\bigr|
  \;=\; \sum_k w_k .
  \label{eq:branching_identity}
\end{equation}
\emph{The sum of the fitted HAR coefficients is the discretised kernel
integral.} Every walk-forward refit therefore already produces an estimate
$\hat{n}_t = \sum_k w_k^{(t)}$ of the branching-ratio analogue at no additional
computational cost, and the sequence $\{\hat{n}_t\}$ over refits is a
measurable reflexivity series. \pending{the $\hat{n}_t$ series and its
regime-conditional and session-conditional decomposition are not yet computed;
the refit coefficient history exists (\texttt{Khist}, 240-bar cadence) and the
estimate must be debiased for the unit ridge used in that fit}

We flag one consistency point in our favour. The reflexivity debate turns on
whether high estimates of $n$ are artifacts of an unmodelled time-varying
background intensity $\mu(t)$: \citet{FilimonovSornette2012} argue that
ignoring intraday seasonality biases $\hat n$ upward toward one. Our Stage-1
diurnal adjustment (Section~\ref{sec:diurnal}) is exactly the $\mu(t)$
correction that this critique demands, adopted originally for an unrelated
reason. Section~\ref{sec:kernel_caveats} argues that it may be too aggressive
in the other direction.

\subsubsection{Two existing negative results become predictions}

The Hawkes reading retro-predicts two of the paper's otherwise awkward
findings.

\paragraph{Bounded nonlinearity.} The quadratic Hawkes model of
\citet{BlancDonierBouchaud2017} makes the feedback a quadratic form in the
past rather than a linear functional, generating volatility clustering and fat
tails endogenously and capturing the Zumbach effect that linear ARCH-type
feedback misses. What it predicts about extractable structure is specific:
real nonlinearity, \emph{additive plus pairwise, and no more}. That is
precisely the saturation reported in Section~\ref{sec:saturation}---the
recoverable nonlinearity is additive-plus-pairwise anchored to the session
clock, and deeper capacity does not pay. A result that reads as a puzzling
ceiling under a generic machine-learning framing reads as a confirmation under
this one.

\paragraph{A low-rank operator.} A multivariate Hawkes process has a matrix
kernel $\Phi_{ij}(u)$, and a low-rank factorisation
$\Phi_{ij}(u) \approx \sum_k a_k\, u^{(k)}_i v^{(k)}_j g_k(u)$ is the natural
structural restriction on it. The operator result---41 exogenous channels
sharing a small number of lag profiles, with quasi-static geometry and fast
amplitude motion---is such a factorisation. Nonparametric estimation of
multivariate Hawkes kernels and their low-rank structure is an established
problem \citep{BacryGaiffasMuzy2015, AchabBacryGaiffas2017}, which gives this
part of the analysis a comparison class it currently lacks.

\subsection{Where the Reading Does Not Yet Hold}
\label{sec:kernel_caveats}

Three discrepancies stand between the correspondence above and a defensible
structural claim. We state them here rather than in a limitations section
because two of them plausibly share a single cause, and that cause is our own
preprocessing.

\paragraph{(i) The fitted kernel takes negative values.} A linear Hawkes
kernel must be non-negative; Equation~\eqref{eq:hawkes_intensity} does not
otherwise define a point process. Our fitted $\varphi$ has negative lobes---
Figure~\ref{fig:target_kernel} plots the positive part. Three readings are
available: a filter artifact (below), genuine inhibition requiring the
nonlinear Hawkes class \citep{BremaudMassoulie1996}, or the concession that
$\varphi$ is a regression weight function and not an intensity kernel at all.
Only the third is safe without further work.

\paragraph{(ii) The exponent is steeper than theory expects.} The
Hawkes-in-finance literature finds kernels near $u^{-1}$---marginally
non-integrable, at the boundary of long memory---and the rough-volatility
correspondence discussed below implies a comparable or shallower decay. Our
$\hat\beta \approx 1.36$ is steeper than both, steep enough that
$\int\varphi$ converges comfortably, i.e.\ short memory in the $L^1$ sense.
This is a substantive disagreement with the literature, not a rounding
difference.

\paragraph{(iii) Both may be artifacts of Stage 1.} Our diurnal adjustment
divides by a \emph{rolling} 20-observation same-slot mean
(Equation~\eqref{eq:diurnal_mean}), which is a high-pass filter rather than a
fixed seasonal profile. High-pass filtering strips low-frequency power, which
biases an estimated decay exponent \emph{steeper}; it also induces negative
side lobes in an estimated kernel. Symptoms (i) and (ii) are exactly the two
signatures of that single cause, and they point in the same direction. The
longest rung of 3{,}125 bars truncates the tail estimate the same way.

\paragraph{The decisive experiment.} Re-estimate the kernel with Stage 1
removed, entering time-of-day as slot dummies in the regression rather than
dividing by a rolling slot mean. If $\hat\beta$ falls toward unity and the
negative lobes vanish, the discrepancies are manufactured by preprocessing and
the structural reading survives; if they persist, the linear Hawkes
correspondence should be abandoned in favour of the weaker
regression-kernel language. \pending{this experiment has not been run; until
it is, we describe the backbone as \emph{power-law-like} and avoid structural
claims about branching ratios}

\subsection{Relation to Rough Volatility}
\label{sec:rough}

Nearly unstable Hawkes processes with heavy-tailed kernels converge, under
rescaling, to rough fractional processes \citep{JaissonRosenbaum2016}, which is
the microstructural route to the rough-volatility models of
\citet{GatheralJaissonRosenbaum2018} and to rough Heston
\citep{ElEuchRosenbaum2019}. In that limit a kernel tail exponent $1+\alpha$
corresponds to a Hurst parameter $H = \alpha - \tfrac{1}{2}$, so the
$H \approx 0.1$ typically reported for volatility corresponds to a tail near
$1.6$. Our $\hat\beta \approx 1.36$ sits closer to that value than to the
$\approx 1$ of the Hawkes-kernel literature. \pending{the exponent conditions
of \citet{JaissonRosenbaum2016} must be checked against $\hat\alpha$ before
this mapping is asserted rather than noted}

This has a direct consequence for the paper's benchmark set that we regard as
a genuine gap: a fractional-kernel forecaster is a \emph{rival in the same
power-law family} as our backbone, and the battery does not contain one. It is
cheap to add, since the direct power-law convolution machinery of
Section~\ref{sec:algorithms} already constructs arbitrary $u^{-\beta}$ features.
\pending{fractional/RFSV arm on the shared panel}

\subsection{The Clock Organ}
\label{sec:clock}

A Hawkes process of the form \eqref{eq:hawkes_intensity} is time-homogeneous.
Intraday realized variance emphatically is not, and the calendar block is the
single largest organ we measure ($-0.0023$ QLIKE, on the organ-ablation cache
rather than the frozen battery). The backbone is therefore
kernel $\times$ clock rather than kernel alone.

The clock does more than scale the kernel. Session-conditional estimation shows
the target kernel \emph{rotating} through the day---cosine similarity to the
pooled kernel of $+0.83$ overnight against $-0.55$ at the open---so the
geometry deforms at session transitions rather than merely breathing in
amplitude. This is the hook that Section~\ref{sec:saturation} pays off: the
nonlinear structure that survives is anchored to exactly these transitions.

\subsection{Summary: The Instrument}
\label{sec:backbone_summary}

The backbone is a truncated power-law self-kernel with a clock organ, of which
the base-5 HAR ladder is a coarse quadrature. Three consequences carry forward:

\begin{enumerate}
  \item \textbf{The incumbent is quadrature-limited.} The base-5 ladder used
    as this paper's shared benchmark is beaten by base-2 OLS by $0.00082$
    QLIKE at DM $t = -10.6$, with no change in information set, model class, or
    tuning. Section~\ref{sec:benchmark_debt} accounts for what this does to
    every increment reported downstream.
  \item \textbf{The natural basis is a mixture of decay exponents.} Direct
    power-law convolution features $K_\beta(t) = \sum_u u^{-\beta} y_{t-u}$,
    normalised, span the same structure as the boxcar ladder but are
    parameterised by decay exponent rather than window length---the natural
    coordinate if the kernel reading is correct, and freshness-robust in a way
    boxcars are not (a 3{,}125-wide rolling mean weights a 65-day-old bar
    exactly as heavily as yesterday's).
  \item \textbf{Two later negative results are predictions, not surprises.}
    Bounded additive-plus-pairwise nonlinearity and a low-rank exogenous
    response operator are what the quadratic and multivariate extensions of
    this backbone respectively predict.
\end{enumerate}
